\subsection{Esiti, costi e unità statistica}
Un successo è un episodio terminato con un circuito compilato, eseguibile sul
Target, e score finito. I controlli del procedimento conservato verificano queste
condizioni prima di pubblicare \texttt{success}. Qualsiasi esito terminale senza
compilazione valida è un fallimento, indipendentemente dalla causa: errore,
timeout o interruzione. I circuiti ancora senza esito sono pendenti e non diventano
fallimenti. Il numero di successi e fallimenti è un conteggio, non una media.

L'unità statistica è il circuito. Se sono conservati più episodi indipendenti
dello stesso circuito e metodo, si calcola prima la media delle misure note del
circuito e poi la media fra circuiti, con un peso uguale per ciascuno. I tentativi
di correzione dello stesso episodio non sono repliche. Lo score del circuito è
la media dei soli episodi riusciti: gli episodi falliti restano nei conteggi,
senza assegnare loro zero. La somma dei costi riguarda invece tutti gli episodi
misurati. Il piano corrente prevede un episodio per circuito; la capacità del
generatore di leggere più episodi non autorizza nuove ripetizioni del Test.

\paragraph{Retry e token.}
Un retry è una chiamata LLM aggiuntiva oltre la prima nello stesso episodio:
\(r=\max(0,N_{\mathrm{chiamate}}-1)\). Il costo in token somma ingresso e uscita
di tutte le chiamate, comprese le correzioni e i tentativi falliti quando
misurabili. L'ingresso è il prompt completo tokenizzato, anche quando il server
riusa una cache; l'uscita è il conteggio restituito dal server. Un token ignoto
non vale zero. Le somme parziali note dei token sono conservate separatamente.
Token, retry e risposta LLM non si applicano a Random e MQT Predictor.

\paragraph{Tempi.}
Il tempo totale decorre dall'ingresso del circuito nella procedura fino all'esito
misurato: include lettura e preparazione, recupero RAG, chiamate, correzioni,
avvio del processo, compilazione e controlli. Il cronometro si ferma prima del
riepilogo finale dei token e della scrittura di \texttt{esito.json}; non include
avvio del server, controlli preliminari di tutta l'esecuzione o generazione di
questo rapporto. La somma di questi tempi è un costo cumulativo, non il tempo
trascorso sul calendario fra due avvii o sessioni.

Il tempo interno di compilazione misura soltanto \texttt{transpile} o
\texttt{rl\_compile}. Il tempo del processo comprende anche importazioni, avvio,
controlli e salvataggi. Se il processo termina per timeout senza pubblicare la
misura interna, quella durata resta mancante: il limite di 100 s non viene
sostituito al tempo del compilatore. La risposta LLM misura l'intervallo fra
invio della generazione e ricezione della risposta completa; con più chiamate
si riporta la loro somma. Si tratta della latenza osservata dal client, non del
solo calcolo dei token. Tutti i tempi sono in secondi, misurati con
\texttt{perf\_counter}.

\paragraph{Copertura.}
Il simbolo -- indica misura mancante o non applicabile; le didascalie distinguono
i casi. Ogni media dichiara quanti circuiti hanno almeno una misura e ogni
riepilogo dichiara gli episodi misurati sul totale pertinente. Una media basata
su misure parziali resta tale; i file leggibili da programma indicano anche
quanti episodi mancano per ogni circuito. Non si eliminano gli esiti sfavorevoli.

\subsection{Confronto di qualità}
La media descrittiva di un metodo usa i suoi successi. Per confrontare due metodi
si usano soltanto i circuiti con uno score disponibile per entrambi, mantenendo
l'appaiamento per identificativo e SHA-256. I fallimenti restano visibili nella
tabella di affidabilità; il confronto sui successi comuni è condizionato alla
riuscita e non li sostituisce. Con episodi ripetuti, un circuito misto entra con
la media dei suoi successi e i fallimenti restano dichiarati.

Per ogni confronto con LLM + RAG si calcola
\(\Delta_c=s_{c,\mathrm{RAG}}-s_{c,\mathrm{altro}}\).
Il piano descrittivo usa 10000 ricampionamenti appaiati con reinserimento dei
circuiti, seed 20260901 e intervallo percentile al 95\% della differenza media.
Con meno di due circuiti comuni l'intervallo non viene stimato. Non si eseguono
test confermativi né si applica il precedente piano sul regret della validation.
L'intervallo descrive la variabilità fra i circuiti osservati: non comprende la
variabilità di nuove esecuzioni, hardware diversi o altre scelte di Random.
Non si dichiara superiorità statistica sulla base di queste sole analisi.
